\section{结论}

\begin{frame}{结论}
  \begin{itemize}
    \item 本文从融资契约视角研究数字化转型如何重塑企业投资结构。
    \item 理论上，融资约束缓解后有形资本调整方向取决于抵押能力与临界阈值的相对位置，研发投入扩张则是无条件的。
    \item 实证上，两化融合贯标显著降低 CapEX1、提高 RD1，推动投资从实物资本向研发活动重配。
    \item 机制上，数字化提升经营信息可验证性，拓宽信贷边界并弱化实物抵押依赖。
    \item 结果上，抵押囤积出清伴随 MRPK 离散度下降和投资对成长机会敏感性上升。
  \end{itemize}
\end{frame}

\begin{frame}{政策含义}
  \begin{itemize}
    \item 推进企业数字化建设时，应将数据治理、流程留痕和经营信息可验证性纳入金融基础设施视角。
    \item 缓解中小企业融资困境不能只扩大抵押物范围，还需要降低银企信息不对称。
    \item 贯标后应配套冗余产能的有序退出通道，避免存量资产流动性折价吞噬数字化红利。
    \item 对实物投资下降的评价应结合研发扩张和配置效率变化，避免简单等同于实体经济空心化。
  \end{itemize}
\end{frame}
